% Generated by scripts/assemble.py; corrections belong in source/corrections.json.
\examhead{57}{南开大学 2026 年数学分析试卷}{57}
% p058-b04; source PDF page 58
\blockanchor{p058-b04}
\question{1}设常数 $\alpha > 0$, 计算积分 $\int_{1}^{+\infty} \frac{\cos(\ln x)}{x^{\alpha + 1}} \, \mathrm{d}x$.\par
% p058-b05; source PDF page 58
\par\Needspace{6\baselineskip}
\blockanchor{p058-b05}
\question{2}设 $L$ 为闭区域 $D = \{(x, y) \mid 0 \leq x \leq 1, x^3 \leq y \leq \sqrt[3]{x}\}$ 的边界, 取逆时针方向, 计算曲线积分\par
\nopagebreak[4]
% p058-b06; source PDF page 58
\blockanchor{p058-b06}
\displayformula{I = \int_ {L} (x ^ {3} y + e ^ {y} + \ln (e ^ {x} + e ^ {y})) \mathrm{d} x + (2 x - 2 y + x y ^ {3} + x e ^ {y}) \mathrm{d} y.}
% p058-b07; source PDF page 58
\par\Needspace{6\baselineskip}
\blockanchor{p058-b07}
\question{3}设 $\Omega$ 为 $\mathbb{R}^n$ 中的有界区域, 函数 $f(x)$ 在 $\overline{\Omega}$ 上连续, 对满足 $U \subseteq \Omega$ 的任意开集 $U$, 象集 $f(U) = \{f(x) | x \in U\}$ 为 $\mathbb{R}^1$ 的开子集, 证明:\par
\nopagebreak[4]
% p058-b08; source PDF page 58
\blockanchor{p058-b08}
\displayformula{\max _ {\boldsymbol {x} \in \overline {{{\Omega}}}} f (\boldsymbol {x}) = \max _ {\boldsymbol {x} \in \partial \Omega} f (\boldsymbol {x}).}
% p058-b09; source PDF page 58
\blockanchor{p058-b09}
\question{4}令 $f(x) = \sum_{n=2}^{\infty} \frac{1}{(n \ln n)^x}$，证明 $f(x)$ 在 $(1, +\infty)$ 上连续可微.\par
% p058-b10; source PDF page 58
\blockanchor{p058-b10}
\question{5}求函数项级数 $\sum_{n=1}^{\infty} \frac{\sin nx}{2^n}$ 和 $\sum_{n=1}^{\infty} \frac{1}{n2^n} \sin^2 \frac{nx}{2}$ 的和函数.\par
% p058-b11; source PDF page 58
\par\Needspace{6\baselineskip}
\blockanchor{p058-b11}
\question{6}设数列  $\{\lambda_{n}\}$  满足  $\lambda_{n}\in\left(\frac{1}{e},+\infty\right), n=1,2,\cdots,$  且  $\varlimsup_{n\to\infty}\lambda_{n}>\frac{1}{e}.$  令\par
\nopagebreak[4]
% p058-b12; source PDF page 58
\blockanchor{p058-b12}
\displayformula{f _ {0} (x) = x, f _ {n} (x) = \lambda_ {n} e ^ {f _ {n - 1} (x)}, n = 1, 2, \dots .}
% p058-b13; source PDF page 58
\blockanchor{p058-b13}
证明：对任意的 $x > 0$ ，都有 $\lim_{n\to \infty}f_n(x) = +\infty .$\par
% p058-b14; source PDF page 58
\blockanchor{p058-b14}
\question{7}设  $a_{n}=\int_{\frac{1}{2}}^{2}\frac{2^{x}}{(1+x^{2})^{n}}\mathrm{d}x,\ n=1,2,\cdots$ ，证明： $\lim_{n\to\infty}\frac{a_{n+1}}{a_{n}}=\frac{4}{5}$\par
% p058-b15; source PDF page 58
\par\Needspace{6\baselineskip}
\blockanchor{p058-b15}
\question{8}证明无穷乘积 $\prod_{n=1}^{\infty}\left(1 + \frac{x}{n^2}\right)$ 对于任意的 $x > 0$ 都收敛, 并求极限\par
\nopagebreak[4]
% p058-b16; source PDF page 58
\blockanchor{p058-b16}
\displayformula{\lim _ {x \rightarrow + \infty} \frac {\ln \left(\prod_ {n = 1} ^ {\infty} \left(1 + \frac {x}{n ^ {2}}\right)\right)}{\sqrt {x}}}
